\documentclass[11pt,a4paper]{article}
\usepackage[utf8]{inputenc}
\usepackage[british]{babel}
\usepackage[a4paper,margin=2.3cm]{geometry}
\usepackage{amsmath,amssymb}
\usepackage{graphicx}
\usepackage{tikz}
\usetikzlibrary{calc,angles,quotes,patterns}
\usepackage[most]{tcolorbox}
\usepackage{enumitem}
\usepackage{array}
\usepackage{titlesec}
\usepackage[hidelinks]{hyperref}

\definecolor{bandgreen}{RGB}{46,139,87}
\definecolor{bandblue}{RGB}{30,100,190}
\definecolor{bandamber}{RGB}{220,140,20}
\definecolor{bandpink}{RGB}{210,70,140}
\definecolor{bandpurple}{RGB}{110,50,160}

\newcounter{prob}
% \prob{label}{title}{source}
\newcommand{\prob}[3]{\par\bigskip\refstepcounter{prob}\label{#1}%
  \noindent\textbf{\theprob.\ #2}\hfill{\small\textit{#3}}\par\nopagebreak\smallskip}
\newcommand{\ans}[2]{\par\smallskip\noindent\textbf{\ref{#1}.}\ #2}

\newcommand{\band}[3]{%
  \clearpage
  \begin{tcolorbox}[colback=#1!12,colframe=#1,boxrule=1pt,arc=2pt]
  {\Large\bfseries\color{#1} #2}\hfill{\large #3}
  \end{tcolorbox}}

\newtcolorbox{diarynote}[1]{colback=black!4,colframe=black!45,boxrule=0.6pt,
  arc=2pt,fonttitle=\bfseries\small,title={#1},fontupper=\small}

\setlength{\parindent}{0pt}
\setlength{\parskip}{3pt}
\definecolor{bandgrey}{RGB}{85,85,85}
\begin{document}

\begin{titlepage}
\centering
\vspace*{3cm}
{\Huge\bfseries The Ladies' Diary\par}
\vspace{0.4cm}
{\LARGE The Gregory years, 1826--1834\par}
\vspace{0.8cm}
{\Large Twenty-two problems from nine issues\par}
\vspace{2cm}
{\large\itshape We admire the talents of our fair correspondent. If she will only make
up her mind to carry her own solutions through, instead of referring to Simpson, Hymer,
\&c.\ she may soon vie with Donna Agnesi, and our own Mrs.\ Somerville.\par}
\vspace{0.3cm}
{\small The editor, on Miss L.~L., 1834\par}
\vfill
{\large Working draft\par}
\end{titlepage}

\section*{A word before the problems}

This booklet fills the last gap in the series. It is made from a bound volume of the Diary
for 1825 to 1834, the 122nd to the 131st almanacs, digitised by Google Books. The new
questions set in it run from 1434 to 1583. Those of 1825 (1434 to 1448) are in the
booklet before this one, which we have checked against the answers printed here, so this
booklet draws on the questions of 1826 to 1834. With it, the
mathematics of \textit{The Ladies' Diary} is now covered from the first questions to the
last. The editor throughout was Olinthus Gregory of the Royal Military Academy, Woolwich.

The Diary had grown up. Its prize questions came from the Continent: M.~Brianchon, of
Brianchon's theorem, set the prize question on tetrahedra in 1830, and a year later
``Investigator'' asked for a simple proof of ``Prof.\ Gauss's new principle of
Mechanics''. Questions were borrowed from Gergonne's \textit{Annales} and Quetelet's
\textit{Correspondance}, and sent from Luxembourg, from Trinity College Dublin and from
New York. Samuel Earnshaw of Sheffield, later of Earnshaw's theorem, set
problem~\ref{v-earnshaw}. Charles Gill, then of Hessle near Hull, set
problems~\ref{v-thp} and~\ref{v-pent}; by 1839 he was a professor in New York, still
writing to the Diary.

\subsection*{The women}

Three women stand out, and they are among the best evidence in the whole series.
\begin{itemize}[leftmargin=*,itemsep=1pt]
\item \textbf{Miss Jane Thoburn of Blyth, ``aged 12 years''}, has the first printed
answer to problem~\ref{v-nine}, in 1829. Her solution is short and correct.
\item \textbf{Miss L.~L.} was answering in quantity by 1834: we count her name against at
least eight questions that year, and her solution to problem~\ref{v-sines} is printed
first. She also set problem~\ref{v-ll}. The editor's note on her, on the title page of
this booklet, compares her with Maria Agnesi and Mary Somerville: high praise, with a
nudge to carry her solutions through for herself.
\item \textbf{Mrs S.~Clay} of Moulton, Lincolnshire, the Sarah Ellis of the previous
booklet, answered at least nine questions in 1828.
\end{itemize}

\subsection*{How to read the labels}

\textit{Q1479} is the Diary's own question number. \textit{After Q1542} means the question
has been retold; the answer at the back is then for the retelling. The bands are those of
the first book in this series; there is nothing young enough for Green or Blue.
Mechanics goes in the last section.

\subsection*{How this edition was made}
This edition was prepared with the assistance of Claude, an AI model made by Anthropic.
Under the editor's direction, Claude worked from Google Books' machine-read text of the
volume. It drafted the retellings, solutions, notes and introduction, recomputed every
answer, and typeset the text. Selection, editorial decisions and final responsibility for
the text are the editor's.

%================ AMBER ================
\band{bandamber}{Amber}{ages 11--14}

\prob{v-nine}{Nine numbers}{Q1479, by J.~B.~Luce}
Find nine numbers in arithmetic progression such that the sum of their squares is
seventeen times the product of the third and the seventh.

\prob{v-breakfast}{The breakfast hour}{Q1464, by Hatton Cummings}
``In summer I sit down to breakfast between 7 and 8 o'clock, at the moment when the hour
and minute hands of my clock make equal angles with the horizon. In winter I sit down
between 8 and 9, when the hands are again equally inclined to the horizon.'' At what
times does he breakfast?

\prob{v-seven}{Dividing by 7 and 13}{Q1512, by Charles Gill}
Give a quick way to tell whether a large number divides by 7 or by 13, without dividing.

\prob{v-earnshaw}{Odd numbers}{Q1528, by Samuel Earnshaw}
Show that $1 + 3 + 5 + 7 + \cdots = (1 - 3 + 5 - 7 + \cdots)^2$, both series being taken
to the same number of terms.

\prob{v-polygons}{The angles of regular polygons}{Q1556, by F.~F.~F., Trinity College,
Dublin}
Show that an angle of a regular polygon of $n$ sides is to an angle of one of $n+1$ sides
as $(1 + 2 + \cdots + (n-1)) - 1$ is to $1 + 2 + \cdots + (n-1)$.

%================ PINK ================
\band{bandpink}{Pink}{ages 14--16}

\prob{v-perimeter}{Perimeter equals area}{Q1465, by Thomas Hindmarsh}
Find a rule for right-angled triangles with whole-number sides whose perimeter is
numerically equal to their area, and give two examples.

\prob{v-cot}{Squares over cotangents}{Q1471, by John Ramsay}
Show that in any triangle the sum of the squares of the sides, divided by the sum of the
cotangents of the angles, is four times the area.

\prob{v-hexagon}{The hexagon of squares}{Q1505, by Petrarch}
Draw a square outwards on each side of a triangle and join the outer corners of
neighbouring squares, making a hexagon. Show that the sum of the squares of the six
sides of the hexagon is four times the sum of the squares of the sides of the triangle.

\prob{v-dodecagon}{The regular dodecagon}{Q1574, by the Rev.~John Sampson}
Show that in a regular dodecagon the distance between two opposite sides equals the side
of the dodecagon plus the side of the square inscribed in the same circle.

\prob{v-palindrome}{A palindrome whose square uses every digit}{Q1576, by the Rev.~John
Sampson}
Find a five-digit number of the form $abcba$ whose square contains all ten digits.

\prob{v-maffett}{Sums that are squares}{after Q1542, by Robert Maffett}
Find three different positive whole numbers such that the sum of any two is a square, and
twice the sum of all three is a fourth power.

%================ PURPLE ================
\band{bandpurple}{Purple}{ages 16--17}

\prob{v-ll}{The circle in the triangle}{Q1577, by Miss L.~L.}
$AB$ is the diameter of a semicircle and $C$ a point on it. Drop the perpendicular $CD$ to
$AB$ and join $AC$. Where should $C$ be for the circle inscribed in triangle $ACD$ to be
as large as possible, and how large is it?

\prob{v-sines}{A hundred sines}{Q1557, by John Baines}
For what value of $a$ is the sum of the hundred terms
$\sin a + \sin 3a + \sin 5a + \cdots + \sin 199a$ greatest?

\prob{v-thp}{Triangular, pentagonal and hexagonal}{Q1468, by Charles Gill}
Find a number greater than 1 that is at once triangular, pentagonal and hexagonal.

\prob{v-cubes}{Odd cubes}{Q1529, by the Rev.~John Sampson}
Show that $1^3 + 3^3 + 5^3 + \cdots$, taken to any number of terms, is always a hexagonal
number. When is it a perfect square?

\prob{v-semicircles}{Semicircles through centres of gravity}{Q1456, by Isaac Newton of Wisbech}
A semicircle stands on a diameter $AB$. A second semicircle, also on the line $AB$ and
touching the first at $B$, passes through the centre of gravity of the first; a third
passes through the centre of gravity of the second, and so on. Find the total area of
all the semicircles.

\prob{v-director}{Right-angled tangents}{Q1504, by Petrarch}
Show that the points from which the two tangents to an ellipse are at right angles lie
on a circle, whose radius is the distance between the ends of the major and minor axes.

\prob{v-chords}{All the chords of a regular polygon}{Q1521, by Petrarch}
Join every pair of corners of a regular polygon of $n$ sides inscribed in a circle of
radius $r$. Show that the sum of all the lines, sides included, is
$nr\cot\frac{\pi}{2n}$.

\prob{v-fermat}{Fermat's theorem}{Q1550, by R.~Vigors}
If $a$ is any whole number and $p$ any prime, $a^p - a$ divides by $p$. Give a simple
proof from first principles.

%================ BEYOND ================
\band{bandgrey}{Beyond A-level}{numbers and motion}

\prob{v-pent}{Two pentagonal numbers}{Q1489, by Charles Gill}
Find two pentagonal numbers whose sum and difference are also pentagonal.

\prob{v-sphere}{Two ways to hang a sphere}{Q1495, by W.~Mawson}
A solid sphere of radius 10 inches swings as a pendulum about a point 20 inches from its
centre. About what other point, inside the sphere, will it swing in the same time?

\prob{v-dog}{The dog and the duck}{Q1565, by Nimrod Knickerbocker, New York}
A sportsman shoots a duck, which falls at the edge of the river exactly opposite him. His
spaniel swims for it, always heading straight at the duck, while the steady stream carries
him downstream. Find the dog's path, in terms of the ratio of the stream's speed to his.

%================ ANSWERS ================
\clearpage
\section*{Answers}

\ans{v-nine}{Any arithmetic progression whose first term is 0, for instance
$0, 1, 2, \dots, 8$: the squares add up to 204, which is $17\times2\times6$. Writing the
terms as $x - 4d, \dots, x + 4d$, the condition becomes $17(x^2 - 4d^2) = 9x^2 + 60d^2$, so
$x = 4d$, and the first term $x - 4d$ is 0. This is Jane Thoburn's solution. P.~M'Kay of Woolwich gave the same progression backwards, 8, 7, \dots,
0.}

\ans{v-breakfast}{At $23\tfrac1{13}$ minutes past 7 and $18\tfrac6{13}$ minutes past 8. The
hands are equally inclined to the horizon when they are mirror images in the line from 12
to 6, so their angles from 12 add up to a whole turn. At $m$ minutes past $h$ the hour hand
is at $30h + \tfrac12m$ degrees and the minute hand at $6m$, so
$30h + 6\tfrac12m = 360$. Matthew Paxton gave this rule, and noted that the hands are also
equally inclined when they are together.}

\ans{v-seven}{Since $1001 = 7\times11\times13$, a thousand leaves remainder $-1$ on division
by each. So split the number into blocks of three digits from the right, and add and
subtract them alternately. The number divides by 7, 11 or 13 exactly when the result does.
For example $123{,}456{,}789$ gives $789 - 456 + 123 = 456$, which divides by neither 7 nor
13, so neither does the number.}

\ans{v-earnshaw}{Both sides are $n^2$. The odd numbers add up to $n^2$, and the alternating
sum pairs off into $-2$'s, giving $-n$ when $n$ is even and $n$ when $n$ is odd.}

\ans{v-polygons}{An angle of a regular $n$-gon is $180^\circ(n-2)/n$, so the ratio is
$\dfrac{(n-2)(n+1)}{n(n-1)}$. And $1 + 2 + \cdots + (n-1) = \tfrac12n(n-1)$, so the
proposed ratio is $\dfrac{\frac12n(n-1) - 1}{\frac12n(n-1)} = \dfrac{(n-2)(n+1)}{n(n-1)}$.}

\ans{v-perimeter}{If the legs are $a$ and $b$, the condition reduces to
$(a - 4)(b - 4) = 8$, so the only whole-number triangles are 5, 12, 13 and 6, 8, 10. Mrs
S.~Clay of Moulton was among the solvers. (Question 1509, two years later, asked for the
two legs alone to add up to the area. That gives $(a-2)(b-2) = 4$, so legs 3 and 6, or 4
and 4, and neither triangle has a whole-number hypotenuse.)}

\ans{v-cot}{$\cot A = (b^2 + c^2 - a^2)/4\Delta$, and similarly for the others. Adding, the
cotangents sum to $(a^2 + b^2 + c^2)/4\Delta$.}

\ans{v-hexagon}{The three sides of the hexagon that are sides of the triangle contribute
the sum $S$ of the squares of the triangle's sides. Each of the other three joins the outer
corners of two squares at a vertex, say $C$. That join is the third side of a triangle with
sides $a$, $b$ and included angle $180^\circ - C$, so its square is
$a^2 + b^2 + 2ab\cos C$. Adding the three and using $2ab\cos C = a^2 + b^2 - c^2$ gives
$3S$, so the whole is $4S$.}

\ans{v-dodecagon}{In a circle of radius $R$ the side of the dodecagon is $2R\sin15^\circ$ and
the side of the inscribed square is $R\sqrt2 = 2R\sin45^\circ$. The distance between
opposite sides is $2R\cos15^\circ$. And $\sin15^\circ + \sin45^\circ =
2\sin30^\circ\cos15^\circ = \cos15^\circ$.}

\ans{v-palindrome}{There are exactly three: $35853^2 = 1{,}285{,}437{,}609$,
$84648^2 = 7{,}165{,}283{,}904$ and $97779^2 = 9{,}560{,}732{,}841$. The Rev.~John Sampson's
printed answer gives all three.}

\ans{v-maffett}{For example 904, 1,400 and 2,696. The pairwise sums are $48^2$, $60^2$ and
$64^2$, and twice the total is $10{,}000 = 10^4$. Writing the pairwise sums as $x^2$, $y^2$,
$z^2$, twice the total is $x^2 + y^2 + z^2$, so one needs three squares adding up to a
fourth power, each smaller than the sum of the other two. Maffett gave a general rule of
this kind. The question allowed equal numbers, and then 72, 72 and 504 also work.}

\ans{v-ll}{When $\cos\angle CAB \approx 0.8269$, that is $\angle CAB \approx 34^\circ13'$; the
largest circle then has diameter about $0.3219$ of $AB$. With $AB = d$ and
$\angle CAB = \phi$, the triangle has $AC = d\cos\phi$, $AD = d\cos^2\phi$ and
$CD = d\cos\phi\sin\phi$, and its inscribed circle has diameter
$AD + CD - AC = d\cos\phi(\cos\phi + \sin\phi - 1)$. Miss L.~L.'s own answer, printed in
1835, gives $34^\circ12'51''$ and $0.321869\,d$, exactly right.}

\ans{v-sines}{About $40'4''$, or $a \approx 0.011656$ radians. The sum is
$\sin^2(100a)/\sin a$, which is greatest where $\tan100a = 200\tan a$. This is the answer of
John Baines, the proposer, and W.~Clough. The solution printed first, by Miss L.~L., gave
$37'47''$; see the grey box.}

\ans{v-thp}{40,755, which is the 285th triangular, the 165th pentagonal and the 143rd
hexagonal number. Every hexagonal number is triangular, so one needs a hexagonal number
that is pentagonal, which leads to a Pell equation. The next is 1,533,776,805. Mr Mason of
Scoulton found 40,755.}

\ans{v-cubes}{The sum of the first $n$ odd cubes is $n^2(2n^2 - 1)$, which is the hexagonal
number $m(2m - 1)$ with $m = n^2$. It is a square when $2n^2 - 1$ is a square, which is
a Pell equation: $n = 1, 5, 29, 169, 985, \dots$, each from the last two by
$n_{k+1} = 6n_k - n_{k-1}$. For $n = 5$ the sum is $1{,}225 = 35^2$.}

\ans{v-semicircles}{About $1.534$ times the first. The centre of gravity of a semicircle of
radius $R$ lies on its axis, $4R/3\pi$ from the diameter. A circle through $B$ with its
centre on $AB$ that passes through that point has radius
$kR$, with $k = \tfrac12\bigl(1 + \tfrac{16}{9\pi^2}\bigr) \approx 0.5901$. The semicircles
shrink by the factor $k$ each time, so their areas add up to $1/(1 - k^2)$ times the first.}

\ans{v-director}{For the ellipse $x^2/a^2 + y^2/b^2 = 1$, the line $y = mx + c$ is a tangent
when $c^2 = a^2m^2 + b^2$. The tangents through $(X, Y)$ have slopes satisfying
$(Y - mX)^2 = a^2m^2 + b^2$, and they are perpendicular when the product of the slopes is
$-1$, which gives $X^2 + Y^2 = a^2 + b^2$. The radius $\sqrt{a^2+b^2}$ is the distance
between the ends of the two axes.}

\ans{v-chords}{The chords from one corner have lengths $2r\sin(k\pi/n)$ for
$k = 1, \dots, n-1$, and $\sum_k\sin(k\pi/n) = \cot(\pi/2n)$. Each chord is counted twice
when we sum over all $n$ corners, so the total is
$\tfrac12n\cdot2r\cot(\pi/2n) = nr\cot\frac\pi{2n}$.}

\ans{v-fermat}{By induction on $a$. The binomial theorem gives
$(a+1)^p - (a+1) = (a^p - a) + \sum_{k=1}^{p-1}\binom pk a^k$, and each $\binom pk$ with
$0 < k < p$ divides by $p$, because $p$ is prime and appears in the numerator but not the
denominator. So if $p$ divides $a^p - a$, it divides $(a+1)^p - (a+1)$.}

\ans{v-pent}{The smallest pair of ordinary pentagonal numbers is 1,560,090 and 7,042,750,
whose sum is 8,602,840 and difference 5,482,660, all four pentagonal. This was found only
in the computer age. John Whitley's printed answer gives pairs such as 7 and 5, and 57 and
35; see the grey box.}

\ans{v-sphere}{2 inches from the centre. Hung at distance $d$ from its centre, the sphere
swings like a simple pendulum of length $d + \tfrac25r^2/d$. At $d = 20$ this is 22 inches.
The same length comes from $d = \tfrac25r^2/20 = 2$: the point of suspension and the centre
of oscillation can always be exchanged.}

\ans{v-dog}{Put the duck at the origin and the dog at $(a, 0)$ to start, on the far bank,
with the stream carrying him in the $y$ direction, and let $k$ be the stream's speed
divided by his. His path is
\[ y = \frac a2\left[\left(\frac xa\right)^{1-k} - \left(\frac xa\right)^{1+k}\right]. \]
If $k < 1$ he reaches the duck. If $k = 1$ he ends at the point $\tfrac12a$ downstream of
it, on its bank, swimming hard and getting nowhere. If $k > 1$ he never reaches the bank
at all. Knickerbocker's own answer, printed in 1834, reaches the same curve in polar
coordinates, and lets the dog start anywhere.}

\begin{diarynote}{When the Diary was wrong}
\textbf{Problem~\ref{v-sines}.} Miss L.~L. expanded $\tan100a$ as a series and cut it off
after two terms, which gives $\tan^2a = 6/(5n^2 - 3n - 2)$ and $37'47''$. The cut is too
early for an angle this large, and the exact condition gives $40'4''$, as Baines and Clough
found. The editor printed her solution first and theirs after it, without
comment.

\textbf{Problem~\ref{v-pent}.} Whitley used the test that $24p + 1$ is a square. That test
also admits the ``pentagonal numbers of negative rank'', $n(3n-1)/2$ with $n$ negative,
such as 2, 7 and 57. With them his pairs work; with ordinary pentagonal numbers they do
not, and the smallest true pair runs into millions.
\end{diarynote}



\vfill
{\small Source: \textit{The Ladies' Diary} for 1825--1834, digitised by Google Books.}

\end{document}
